
\documentclass[12pt]{article}
\usepackage{amsmath}
\usepackage{amsfonts}
\usepackage{amssymb}
\usepackage{geometry}
\usepackage{graphicx}
\usepackage{accents}
\geometry{a4paper, margin=0.8in}
\pagenumbering{gobble}

\title{Homework assignment for 02YELMA \\ Chapter: Electromagnetic waves}
\begin{document}
\date{}
\maketitle
\vspace{-0.8in}

\begin{enumerate}

    \item We know that if electric fields $\vec{E}$ and magnetic fields $\vec{B}$ exist in vacuum (in absence of charges and currents), they behave like waves according to Maxwell's equations. For now, let us focus only on electric fields. Consider that $\vec{E}$ propagates in the $z$ direction and behaves like plane waves given in the form $\accentset{\sim}{\vec{E}}(z,t)=\accentset{\sim}{\vec{E_0}}e^{i(kz-\omega t)}$ where the phase $\delta$ is absorbed in the constant $\accentset{\sim}{\vec{E_0}}\equiv \vec{E_0}e^{i\delta}$ and $\vec{E_0}$ is the real amplitude of the electric field. Note that the real amplitude can be written in general (in 3 dimensions) as $\vec{E_0}= (E_0)_x\hat{x}+(E_0)_y\hat{y}+(E_0)_z\hat{z}$. By considering $\vec{\nabla}.\vec{E}=0$ (or equivalently, $\vec{\nabla}.\accentset{\sim}{\vec{E}}=0)$, show that the component of electric field in the propogation direction is zero, i.e., $(E_0)_z=0$. Is this also valid for magnetic fields? 

    \item Consider Faraday's law $\vec{\nabla}\times\vec{E}=-\frac{\partial{\vec{B}}}{\partial{t}}$ and $\accentset{\sim}{\vec{E}}(z,t)=\accentset{\sim}{\vec{E_0}}e^{i(kz-\omega t)}$ to show that $(\accentset{\sim}{B_0})_y(\omega/k)=(\accentset{\sim}{E_0})_x$ and $(\accentset{\sim}{B_0})_x(\omega/k)=-(\accentset{\sim}{E_0})_y$. If we also consider $(E_0)_z=(B_0)_z=~0$, explain how we can summarize these four equations with a single equation:
    \begin{equation*}
    (\accentset{\sim}{B_0})=\frac{k}{\omega}(\hat{z}\times \accentset{\sim}{\vec{E}}_0).
    \end{equation*}
    Using the equation above as guidance, briefly describe how electric and magnetic fields behave in vacuum.



    \item A plane electromagnetic wave in vacuum propagates in the $+z$-direction. Its electric field is given by

\[
\vec{E}(z,t)=E_0\cos(kz-\omega t)\,\hat{x}
\]

where $E_0$, $k$, and $\omega$ are positive constants.

\begin{enumerate}

\item Using Maxwell’s equations, determine the corresponding magnetic field $\vec{B}(z,t)$.

\item Compute the instantaneous Poynting vector

\[
\vec{S}=\frac{1}{\mu_0}\vec{E}\times\vec{B}
\]

and show explicitly that it points in the direction of propagation.

\item Find the time-averaged Poynting vector $\langle \vec{S} \rangle$.

\item Show that the instantaneous electromagnetic energy density

\[
u=\frac{1}{2}\varepsilon_0 E^2+\frac{1}{2\mu_0}B^2
\]

can be written entirely in terms of the electric field amplitude $E_0$, and prove that the electric and magnetic fields contribute equally to the total energy density.


\end{enumerate}

\item A parallel-plate capacitor with circular plates of radius $R$ is being charged by a steady current $I$. The plate separation is $d$, with $d \ll R$, so edge effects may be neglected.

At a given instant, the electric field between the plates is uniform and directed along the $z$-axis:

\[
\vec{E}(t)=E(t)\,\hat{z}.
\]

Although no conduction current passes through the vacuum gap between the plates, energy is still transferred into the capacitor.

\begin{enumerate}
    \item Using Maxwell–Ampère’s law, determine the magnetic field $\vec{B}(r)$ between the plates for $r<R$, where $r$ is the distance from the central axis.

    \item Compute the Poynting vector
    \[
    \vec{S}
    =
    \frac{1}{\mu_0}\vec{E}\times\vec{B}
    \]
    in the region between the plates, and determine its direction.

    \item Show that electromagnetic energy flows into the capacitor from the sides rather than along the wire axis.

    \item Calculate the total rate of energy flow into the volume between the plates by integrating the Poynting vector over the cylindrical side surface.

    \item Show that this equals the rate of increase of electrostatic energy stored in the capacitor:
    \[
    U
    =
    \frac{1}{2}\varepsilon_0 E^2(\pi R^2 d).
    \]

    \item Explain conceptually why this example is important for understanding displacement current and the physical meaning of the Poynting vector.

\end{enumerate}

\item A plane electromagnetic wave propagates in vacuum in the \(+z\)-direction. Its electric field is given by

\[
\vec{E}(z,t)
=
\hat{\vec{x}}\,E_0 \cos(kz-\omega t)
+
\hat{\vec{y}}\,E_0 \cos\!\left(kz-\omega t+\phi\right),
\]

where \(E_0\) is a real positive constant and \(\phi\) is a constant phase difference.

\begin{enumerate}
    \item Eliminate the time dependence to obtain the trajectory traced by the tip of the electric field vector in the \(xy\)-plane at a fixed position \(z=z_0\).

    \item Determine the special values of \(\phi\) for which the wave is:
    \begin{itemize}
        \item linearly polarized,
        \item circularly polarized.
    \end{itemize}

    \item For the circular polarization case, determine whether the rotation is clockwise or counterclockwise when viewed looking in the \(+z\)-direction.

    \item Using the corresponding magnetic field, verify that the time-averaged Poynting vector points along the direction of propagation and explain why its magnitude is independent of \(\phi\) even though the polarization state changes.
\end{enumerate}


\end{enumerate}





\end{document}
